\documentclass[11pt,a4paper]{article}

\usepackage[utf8]{inputenc}
\usepackage[T1]{fontenc}
\usepackage[french]{babel}
\usepackage[a4paper,margin=2.2cm]{geometry}
\usepackage{amsmath,amssymb,amsthm,mathtools}
\usepackage{lmodern}
\usepackage{enumitem}
\usepackage{hyperref}
\usepackage{fancyhdr}
\usepackage{xcolor}
\usepackage{stmaryrd}

\hypersetup{
    colorlinks=true,
    linkcolor=black,
    urlcolor=blue
}

\setlength{\parindent}{0pt}
\setlength{\parskip}{0.6em}

\pagestyle{fancy}
\fancyhf{}
\fancyhead[L]{Fiche d’exercices --- Algèbre linéaire}
\fancyhead[R]{Chapitre 1 : Ensembles et raisonnement}
\fancyfoot[C]{\thepage}

\newcommand{\R}{\mathbb{R}}
\newcommand{\N}{\mathbb{N}}
\newcommand{\Z}{\mathbb{Z}}
\newcommand{\Q}{\mathbb{Q}}
\newcommand{\C}{\mathbb{C}}
\newcommand{\PP}{\mathcal{P}}

\begin{document}

\begin{center}
    {\LARGE \textbf{Fiche d’exercices --- Chapitre 1}}\\[0.3em]
    {\Large \textbf{Vocabulaire ensembliste et méthodes de raisonnement}}\\[0.5em]
    {\large Algèbre linéaire}
\end{center}

\hrule
\vspace{0.8em}

AIDE de lecture type d’exo :  
\textbf{APP} = application directe, \textbf{LOG} = logique et quantificateurs, \textbf{DEM} = démonstration, \textbf{PREP} = réflexion plus poussée.  
Les étoiles représentent un gradient de difficulté de 1 à 5.

\tableofcontents

\newpage

\section*{Exercices d’application directe}
\addcontentsline{toc}{section}{Exercices d’application directe}

\textbf{Exercice 1 [APP \(\star\)]}\\
Appartenance ou non-appartenance

Dire pour chacune des affirmations suivantes si elle est vraie ou fausse :
\begin{enumerate}[label=\alph*)]
    \item \(2\in \N\)
    \item \(-3\in \N\)
    \item \(\pi\in \Q\)
    \item \sqrt{2}\in \R
    \item \(i\in \C\)
    \item \(\frac{3}{4}\in \Z\)
\end{enumerate}

\textbf{Exercice 2 [APP \(\star\)]}\\
Inclusions d’ensembles

Dire si les inclusions suivantes sont vraies :
\begin{enumerate}[label=\alph*)]
    \item \(\N\subset \Z\)
    \item \(\Z\subset \Q\)
    \item \(\Q\subset \R\)
    \item \(\R\subset \C\)
    \item \(\R\subset \Q\)
\end{enumerate}

\textbf{Exercice 3 [APP \(\star\star\)]}\\
Définition par extension et compréhension

Écrire chacun des ensembles suivants :
\begin{enumerate}[label=\alph*)]
    \item par extension lorsque c’est possible ;
    \item par compréhension.
\end{enumerate}

\begin{enumerate}[label=\arabic*)]
    \item l’ensemble des entiers naturels inférieurs ou égaux à \(4\),
    \item l’ensemble des entiers relatifs pairs,
    \item l’ensemble des réels strictement positifs,
    \item l’ensemble des solutions réelles de \(x^2=4\).
\end{enumerate}

\textbf{Exercice 4 [APP \(\star\star\)]}\\
Union, intersection, différence

Soient
\[
A=\{1,2,3,4\},\qquad B=\{3,4,5,6\},\qquad C=\{2,4,6\}.
\]
Déterminer :
\begin{enumerate}[label=\alph*)]
    \item \(A\cup B\),
    \item \(A\cap B\),
    \item \(A\setminus B\),
    \item \(B\setminus A\),
    \item \((A\cup B)\cap C\),
    \item \(A\cap(B\setminus C)\).
\end{enumerate}

\textbf{Exercice 5 [APP \(\star\star\)]}\\
Intervalles et opérations ensemblistes

Déterminer :
\begin{enumerate}[label=\alph*)]
    \item \(]-\infty,2]\cap [1,+\infty[\),
    \item \(]-\infty,2]\cup [1,+\infty[\),
    \item \([0,3]\setminus ]1,2[\),
    \item \(\R\setminus [0,1]\),
    \item \(]-\infty,0[\cap [0,+\infty[\).
\end{enumerate}

\textbf{Exercice 6 [APP \(\star\star\)]}\\
Disjonction

Parmi les couples d’ensembles suivants, préciser lesquels sont disjoints :
\begin{enumerate}[label=\alph*)]
    \item \(]-\infty,0[\) et \([0,+\infty[\),
    \item \([1,3]\) et \([3,5]\),
    \item \(\{1,2\}\) et \(\{3,4\}\),
    \item \(2\Z\) et \(2\Z+1\).
\end{enumerate}

\textbf{Exercice 7 [APP \(\star\star\)]}\\
Produit cartésien

Soient
\[
A=\{1,2\},\qquad B=\{a,b,c\}.
\]
Déterminer :
\begin{enumerate}[label=\alph*)]
    \item \(A\times B\),
    \item \(B\times A\),
    \item le nombre d’éléments de \(A\times B\),
    \item si \(A\times B=B\times A\).
\end{enumerate}

\textbf{Exercice 8 [APP \(\star\star\)]}\\
Égalité de deux ensembles

Montrer les égalités suivantes :
\begin{enumerate}[label=\alph*)]
    \item \(\{x\in \R\mid x^2=9\}=\{-3,3\}\),
    \item \(\{x\in \R\mid x^2-5x+6=0\}=\{2,3\}\),
    \item \(\{n\in \Z\mid n \text{ est multiple de } 2\}=\{2k\mid k\in \Z\}\).
\end{enumerate}

\section*{Exercices de logique et de quantificateurs}
\addcontentsline{toc}{section}{Exercices de logique et de quantificateurs}

\textbf{Exercice 9 [LOG \(\star\star\)]}\\
Traduire avec des quantificateurs

Écrire avec les symboles \(\forall\), \(\exists\), \(\exists!\) les propositions suivantes :
\begin{enumerate}[label=\alph*)]
    \item Tout réel a un carré positif ou nul.
    \item Il existe un réel dont le carré vaut \(2\).
    \item Il existe un unique réel \(x\) tel que \(x+1=0\).
    \item Tout entier naturel possède un successeur.
\end{enumerate}

\textbf{Exercice 10 [LOG \(\star\star\)]}\\
Négation de propositions quantifiées

Nier correctement les propositions suivantes :
\begin{enumerate}[label=\alph*)]
    \item \(\forall x\in\R,\ x^2+1>0\),
    \item \(\exists x\in\R,\ x^2=-1\),
    \item \(\forall n\in\N,\ n+1\in\N\),
    \item \(\exists x\in\R,\ x>3 \text{ et } x<4\).
\end{enumerate}

\textbf{Exercice 11 [LOG \(\star\star\)]}\\
Existence et unicité

Pour chacune des affirmations suivantes, dire si elle est vraie :
\begin{enumerate}[label=\alph*)]
    \item \(\exists! x\in \R,\ x+5=0\),
    \item \(\exists! x\in \R,\ x^2=1\),
    \item \(\exists x\in \R,\ x^2+1=0\),
    \item \(\exists! x\in \R,\ 2x-8=0\).
\end{enumerate}

\textbf{Exercice 12 [LOG \(\star\star\star\)]}\\
Conditions nécessaires et suffisantes

Pour chacune des propositions suivantes, préciser si elle exprime :
\begin{itemize}
    \item une implication,
    \item une réciproque,
    \item une équivalence.
\end{itemize}

\begin{enumerate}[label=\alph*)]
    \item Si \(x=0\), alors \(x^2=0\).
    \item Si \(x^2=0\), alors \(x=0\).
    \item \(x^2=0\) si et seulement si \(x=0\).
    \item Si un entier est multiple de \(4\), alors il est pair.
\end{enumerate}

\textbf{Exercice 13 [LOG \(\star\star\star\)]}\\
Contraposée

Écrire la contraposée de chacune des implications suivantes :
\begin{enumerate}[label=\alph*)]
    \item Si \(x>2\), alors \(x^2>4\).
    \item Si \(n\) est multiple de \(4\), alors \(n\) est pair.
    \item Si \(x^2\neq 0\), alors \(x\neq 0\).
    \item Si \(x\in ]0,1[\), alors \(x>0\).
\end{enumerate}

\textbf{Exercice 14 [LOG \(\star\star\star\)]}\\
Réciproque vraie ou fausse

Pour chacune des implications suivantes :
\begin{enumerate}[label=\alph*)]
    \item dire si elle est vraie ;
    \item dire si sa réciproque est vraie ;
    \item si la réciproque est fausse, donner un contre-exemple.
\end{enumerate}

\begin{enumerate}[label=\arabic*)]
    \item Si \(x=1\), alors \(x^2=1\).
    \item Si \(n\) est multiple de \(6\), alors \(n\) est multiple de \(3\).
    \item Si \(x>1\), alors \(x^2>1\).
    \item Si \(x^2>1\), alors \(x>1\).
\end{enumerate}

\section*{Exercices de démonstration}
\addcontentsline{toc}{section}{Exercices de démonstration}

\textbf{Exercice 15 [DEM \(\star\star\)]}\\
Double inclusion

Montrer que pour tous ensembles \(A\) et \(B\),
\[
A=B \Longleftrightarrow (A\subset B \text{ et } B\subset A).
\]

\textbf{Exercice 16 [DEM \(\star\star\)]}\\
Commutativité de l’union et de l’intersection

Montrer que pour tous ensembles \(A\) et \(B\),
\[
A\cup B=B\cup A
\qquad \text{et} \qquad
A\cap B=B\cap A.
\]

\textbf{Exercice 17 [DEM \(\star\star\)]}\\
Associativité

Montrer que pour tous ensembles \(A\), \(B\), \(C\),
\[
(A\cup B)\cup C=A\cup(B\cup C).
\]

\textbf{Exercice 18 [DEM \(\star\star\star\)]}\\
Distributivité

Montrer que pour tous ensembles \(A\), \(B\), \(C\),
\[
A\cap(B\cup C)=(A\cap B)\cup(A\cap C).
\]

\textbf{Exercice 19 [DEM \(\star\star\star\)]}\\
Formules de De Morgan

Soit \(E\) un ensemble de référence et \(A,B\subset E\). Montrer que
\[
E\setminus(A\cup B)=(E\setminus A)\cap(E\setminus B).
\]

Puis démontrer la seconde formule de De Morgan :
\[
E\setminus(A\cap B)=(E\setminus A)\cup(E\setminus B).
\]

\textbf{Exercice 20 [DEM \(\star\star\)]}\\
Produit nul dans \(\Z\)

Montrer que si \(ab=0\) avec \(a,b\in\Z\), alors \(a=0\) ou \(b=0\).

\textbf{Exercice 21 [DEM \(\star\star\)]}\\
Raisonnement direct

Montrer que la somme de trois entiers pairs est paire.

\textbf{Exercice 22 [DEM \(\star\star\)]}\\
Disjonction de cas

Montrer que pour tout entier \(n\), le produit \(n(n+1)\) est pair.

\textbf{Exercice 23 [DEM \(\star\star\star\)]}\\
Contraposée

Montrer que si \(n^2\) est pair, alors \(n\) est pair.

\textbf{Exercice 24 [DEM \(\star\star\star\)]}\\
Contraposée plus fine

Montrer que si \(n^2\) est multiple de \(3\), alors \(n\) est multiple de \(3\).

\textbf{Exercice 25 [DEM \(\star\star\star\)]}\\
Raisonnement par l’absurde

Montrer par l’absurde que l’équation
\[
x^2=-1
\]
n’admet aucune solution réelle.

\textbf{Exercice 26 [DEM \(\star\star\star\)]}\\
Irrationalité

Montrer que \(\sqrt{2}\notin \Q\).

\textbf{Exercice 27 [DEM \(\star\star\star\)]}\\
Existence et unicité

Montrer que pour tout réel \(a\), l’équation
\[
x+a=0
\]
admet une unique solution réelle.

\textbf{Exercice 28 [DEM \(\star\star\star\)]}\\
Existence et unicité affine

Montrer que pour tout réel \(a\), il existe un unique réel \(x\) tel que
\[
2x+a=0.
\]

\textbf{Exercice 29 [DEM \(\star\star\star\)]}\\
Récurrence simple

Montrer par récurrence que pour tout \(n\in\N\),
\[
1+2+\cdots+n=\frac{n(n+1)}{2}.
\]

\textbf{Exercice 30 [DEM \(\star\star\star\)]}\\
Récurrence sur les entiers impairs

Montrer par récurrence que pour tout \(n\in\N\),
\[
1+3+5+\cdots+(2n+1)=(n+1)^2.
\]

\textbf{Exercice 31 [DEM \(\star\star\star\star\)]}\\
Inégalité classique

Montrer que pour tout \(n\in\N\),
\[
2^n\geq n+1.
\]

\section*{Exercices de réflexion}
\addcontentsline{toc}{section}{Exercices de réflexion}

\textbf{Exercice 32 [PREP \(\star\star\star\)]}\\
Analyse-synthèse

Déterminer tous les réels \(x\) tels que
\[
x+\frac{1}{x}=2.
\]
On demandera une rédaction par analyse-synthèse.

\textbf{Exercice 33 [PREP \(\star\star\star\)]}\\
Égalité d’ensembles un peu plus subtile

Montrer que
\[
A\cap(B\setminus C)=(A\cap B)\setminus C.
\]

\textbf{Exercice 34 [PREP \(\star\star\star\star\)]}\\
Union de deux sous-ensembles

Soient \(A\) et \(B\) deux ensembles. Montrer que
\[
A\subset B \Longleftrightarrow A\cup B=B.
\]

Puis montrer que
\[
A\subset B \Longleftrightarrow A\cap B=A.
\]

\textbf{Exercice 35 [PREP \(\star\star\star\star\)]}\\
Caractérisation de l’ensemble vide

Soit \(A\) un ensemble. Montrer que
\[
A=\varnothing \Longleftrightarrow \forall x,\ x\notin A.
\]

\textbf{Exercice 36 [PREP \(\star\star\star\star\)]}\\
Récurrence avec paramètre

Montrer que pour tout réel \(a\neq 1\) et tout entier \(n\in\N\),
\[
1+a+a^2+\cdots+a^n=\frac{1-a^{n+1}}{1-a}.
\]

\textbf{Exercice 37 [PREP \(\star\star\star\star\)]}\\
Existence et unicité géométrique

Soient \(u\) et \(v\) deux vecteurs d’un espace vectoriel. Montrer qu’il existe un unique vecteur \(x\) tel que
\[
x+u=v.
\]

\textbf{Exercice 38 [PREP \(\star\star\star\star\star\)]}\\
Quantificateurs et négations

Nier correctement la proposition suivante :
\[
\forall x\in\R,\ \exists y\in\R,\ x+y=0.
\]
Puis dire si la proposition initiale est vraie ou fausse.

\textbf{Exercice 39 [PREP \(\star\star\star\star\star\)]}\\
Équivalence logique

Montrer que les assertions suivantes sont équivalentes :
\begin{enumerate}[label=\alph*)]
    \item \(A\subset B\),
    \item \(A\cap (\R\setminus B)=\varnothing\).
\end{enumerate}

\textbf{Exercice 40 [PREP \(\star\star\star\star\star\)]}\\
Synthèse générale

Pour chacune des affirmations suivantes :
\begin{enumerate}[label=\alph*)]
    \item dire si elle est vraie ou fausse ;
    \item si elle est vraie, la démontrer ;
    \item si elle est fausse, donner un contre-exemple.
\end{enumerate}

\begin{enumerate}[label=\arabic*)]
    \item Si \(A\subset B\), alors \(A\cap C\subset B\cap C\).
    \item Si \(A\cup C=B\cup C\), alors \(A=B\).
    \item Si \(A\cap C=B\cap C\), alors \(A=B\).
    \item Si \(A\subset B\) et \(B\subset C\), alors \(A\subset C\).
\end{enumerate}

\section*{Pour aller plus loin}
\addcontentsline{toc}{section}{Pour aller plus loin}

\textbf{Exercice 41 [PREP \(\star\star\star\star\star\)]}\\
Structure logique d’un énoncé

Écrire la négation précise de chacune des propositions suivantes :
\begin{enumerate}[label=\alph*)]
    \item \(\forall x\in\R,\ \exists ! y\in\R,\ x+y=0\),
    \item \(\exists x\in\R,\ \forall y\in\R,\ x+y=0\),
    \item \(\forall n\in\N,\ \exists p\in\N,\ p>n\).
\end{enumerate}

\textbf{Exercice 42 [PREP \(\star\star\star\star\star\)]}\\
Analyse-synthèse et unicité

Déterminer tous les réels \(x\) tels que
\[
x^2-2x+1=0
\]
par une rédaction complète en analyse-synthèse, puis expliquer pourquoi cette méthode apporte plus qu’une simple résolution calculatoire.

\textbf{Exercice 43 [DEM \(\star\star\star\star\)]}\\
Principe de récurrence renforcé

Montrer par récurrence que, pour tout \(n\in\N\),
\[
\sum_{k=0}^n (2k+1)=(n+1)^2.
\]

\textbf{Exercice 44 [DEM \(\star\star\star\star\)]}\\
Preuve en deux sens

Montrer que pour tout réel \(x\),
\[
x^2=1 \Longleftrightarrow x=1 \text{ ou } x=-1.
\]

\textbf{Exercice 45 [PREP \(\star\star\star\star\star\)]}\\
Problème de synthèse final

Soient \(A,B,C\) trois ensembles. Étudier la validité des implications suivantes :
\begin{enumerate}[label=\alph*)]
    \item \(A\subset B \Longrightarrow A\cup C\subset B\cup C\),
    \item \(A\subset B \Longrightarrow A\setminus C\subset B\setminus C\),
    \item \(A\subset B \Longrightarrow C\setminus B\subset C\setminus A\).
\end{enumerate}
On justifiera rigoureusement chaque réponse.

\end{document}